\chapter{ИССЛЕДОВАНИЕ: ЭКСПЕРИМЕНТЫ И ВЫБОР ТЕХНОЛОГИЙ}

Глава содержит экспериментальную постановку задачи, обоснование
технологического стека и проверенные результаты текущего MVP:
собранный корпус данных, обзор практик решения задачи и выбранную
архитектуру.

% ─────────────────────────────────────────────────────────────────
\section{Дизайн эксперимента}
% ─────────────────────────────────────────────────────────────────

\subsection{Тестовый корпус и оценочные артефакты}

Для основной линии исследования зафиксирован
\textbf{измеренный корпус одношаговой LLM-конфигурации (B2) из 19 документов},
репрезентирующих три категории источников, с~которыми должен работать
ИИ-методист в~B2B- и образовательном контуре:

\begin{itemize}
  \item \textbf{7 продуктово-редакционных документов:}
    краткое описание продукта, брендовые и редакционные руководства,
    политики стиля и HR-коммуникаций;
  \item \textbf{6 корпоративных регламентов:}
    политики информационной безопасности, персональных данных,
    управления инцидентами, адаптации сотрудников и антикоррупционного поведения;
  \item \textbf{6 открытых учебных программ MIT~OCW:}
    англоязычные учебные планы по экономике, математике, программированию
    и физике~\cite{mit_ocw_2024}.
\end{itemize}

В измеряемое ядро включались документы, пригодные для технического
базового прогона без раскрытия закрытых внешних данных. Для каждого
документа в артефакте базового сравнения сохранены
идентификатор, категория, язык, время генерации и рассчитанные метрики.

Корпус охватывает три существенно разных типа входа: редакционные
правила, регламентные тексты и академические программы. Это соответствует
целевым сценариям продукта: подготовка внутреннего курса по продукту,
адаптационный курс по правилам компании и асинхронный учебный курс.

Для автоматизированной оценки использованы также два вспомогательных
оценочных набора:

\begin{itemize}
  \item \textbf{30 вопросно-ответных пар} (\texttt{qa\_pairs.json}) для
    детерминированной оценки по схеме RAG на уровне ответов на вопросы;
  \item \textbf{100 утверждений} (\texttt{statements.json}), из которых
    50 корректных и 50 намеренно искажённых, для оценки проверки фактов.
\end{itemize}

\subsection[Конфигурации эксперимента и расшифровка B1--B4]{Конфигурации эксперимента: ручной подход (B1), одношаговая LLM (B2), ИИ-конструктор (B3) и ИИ-методист (B4)}

Для исследования определены четыре конфигурации. В качестве количественного
доказательства текущего MVP используются одношаговая LLM-конфигурация (B2)
и ИИ-методист (B4), а ручная разработка (B1) и ИИ-конструктор (B3) сохранены как
контрольные линии для расширения методики:

\begin{itemize}
  \item \textbf{B1} --- ручная разработка структуры курса
    методистом-экспертом; используется как теоретическая верхняя граница
    качества, но не заявляется как выполненный замер;
  \item \textbf{B2} --- одношаговая генерация на базовой LLM без агентской
    оркестрации; это базовая конфигурация, полностью прогнанная на ядровом корпусе;
  \item \textbf{B3} --- коммерческий инструмент класса конструкторов курсов на~основе ИИ
    (в работе закладывается Coursebox AI как ближайший рыночный аналог);
    используется для сравнительного анализа рынка, но не как измеренная
    экспериментальная линия;
  \item \textbf{B4} --- ИИ-методист: мультиагентный конвейер поддержки
    методиста, объединяющий обработку документов, поиск релевантных фрагментов,
    проектирование курса, подготовку черновых материалов, проверку
    фактов и~проверка.
\end{itemize}

Статус базовых конфигураций на момент написания:

\begin{itemize}
  \item \textbf{B2} дал воспроизводимые количественные результаты
    по 19 документам и служит базовым ориентиром;
  \item \textbf{B4} завершил парный прогон на тех же 19 документах
    и используется как основной измеренный результат ИИ-методиста;
  \item \textbf{B1} и \textbf{B3} не используются как количественное
    доказательство гипотез H1--H3 в текущей работе.
\end{itemize}

Глава~2 фиксирует измеренный экспериментальный слой одношаговой LLM (B2)
и ИИ-методиста (B4) и
дополнительную автоматизированную проверку RAG и фактологическую проверку; границы
интерпретации этих результатов отдельно указаны в разделе ограничений.
Отдельно подготовлено описание корпуса для следующего шага:
сравнительной проверки на открытых курсах $n \geq 30$. Она нужна именно
для внешнего сравнительного утверждения, тогда как текущие 19 документов
остаются лабораторной проверкой MVP текущей инженерной версии.

\subsection{Метрики оценки}

На текущем этапе используются четыре группы метрик:

\begin{itemize}
  \item \textbf{структурные}: валидность JSON-структуры, число модулей,
    число уроков, наличие обязательных полей;
  \item \textbf{педагогические}: покрытие уровней таксономии Блума --- доля из шести
    уровней таксономии Блума, представленных в учебных целях курса;
  \item \textbf{поисково-фактологические}: достоверность относительно
    источника (faithfulness), релевантность ответа (answer relevancy),
    точность контекста (context precision) и метрики проверки фактов;
  \item \textbf{операционные}: время подготовки черновика, плотность ссылок на~источник,
    себестоимость курса, пригодность к локальному развёртыванию.
\end{itemize}

Набор охватывает форму результата, педагогическую глубину, связь
с источником и эксплуатационную стоимость.

Вместе с тем одной таксономии Блума недостаточно для полноценной
педагогической и пользовательской оценки курса. Она хорошо описывает
когнитивный уровень целей, но слабее отражает структурную сложность
результата, согласованность курса, качество онлайн-дизайна и эффект
для обучающегося. Поэтому в~методологическом контуре исследования
рассматривается более широкий набор рамок
~\cite{monash_solo_2026, webb_dok_guide_2013,
advancehe_constructive_alignment_2026, quality_matters_he_rubric_2026,
fink_significant_learning_2026, chi_wylie_2014, cast_udl_2024}.

\subsection{Расширенная педагогическая и пользовательская рамка}

\begin{table}[h!]
  \centering
  \small
  \caption{Дополнительные рамки оценки качества курса сверх таксономии Блума}
  \label{tab:extended-learning-metrics}
  \begin{tabular}{|p{3.2cm}|p{4.0cm}|p{5.3cm}|}
  \hline
  \textbf{Рамка} & \textbf{Что измеряет} & \textbf{Потенциальная операционализация для ИИ-методиста} \\
  \hline
  Таксономия SOLO
    & структурную сложность и~степень интеграции результата обучения
    & экспертная оценка структуры курса по~уровням
      от одноструктурного до расширенного абстрактного \\
  \hline
  Глубина знания по Webb (DOK)
    & когнитивную строгость заданий и~оценочных материалов
    & распределение сгенерированных заданий по~уровням DOK~1--4 \\
  \hline
  Конструктивное согласование (Biggs)
    & согласованность целей, учебных активностей и~оценивания
    & доля целей, для~которых в~курсе есть связанная активность
      и~способ проверки \\
  \hline
  Методика Quality Matters
    & качество онлайн-курса как цифрового образовательного продукта
    & итоговый балл по~целям, оцениванию, материалам,
      активностям, технологиям и~доступности \\
  \hline
  Таксономия значимого обучения Fink
    & значимое обучение сверх простого воспроизведения содержания
    & покрытие компонентов применения, интеграции,
      личностного измерения, ценностного отношения и умения учиться \\
  \hline
  Рамка ICAP
    & тип учебной активности и~глубину вовлечения обучающегося
    & доля пассивных, активных, конструктивных и интерактивных
      учебных активностей \\
  \hline
  UDL: Universal Design for Learning
    & доступность, вариативность подачи и~варианты действия
    & покрытие по~осям engagement, representation,
      action \& expression \\
  \hline
\end{tabular}
\end{table}

\begin{table}[h!]
  \centering
  \small
  \caption{Показатели эффекта обучения, ориентированные на обучающегося}
  \label{tab:learner-centered-metrics}
  \begin{tabular}{|p{3.2cm}|p{4.0cm}|p{5.3cm}|}
  \hline
  \textbf{Показатель} & \textbf{Что измеряет} & \textbf{Потенциальная операционализация} \\
  \hline
  Самоэффективность (Self-efficacy)
    & уверенность обучающегося, что он может выполнить конкретную задачу
    & Likert-вопросы, ориентированные на~конкретную задачу,
      вида
      «теперь я могу самостоятельно сделать X» \\
  \hline
  Ощущаемая компетентность (Perceived Competence)
    & субъективное ощущение владения темой или навыком
    & короткая посткурсовая шкала PCS по~конкретному домену \\
  \hline
  Перенос обучения в~практику (Transfer of learning)
    & факт применения изученного в~рабочей или учебной ситуации
    & отложенный опрос через 2--4~недели: применял~ли, сколько раз,
      какие барьеры возникли \\
  \hline
  Retrospective pre-post
    & изменение восприятия своих возможностей с~учётом эффекта
      response-shift
    & один пост-опросник с~оценкой состояния
      «до курса» и~«после курса» \\
  \hline
  Достижение цели (Goal attainment)
    & степень достижения индивидуальной учебной или рабочей цели
    & мини-шкала достижения цели по~цели, сформулированной
      перед~обучением \\
  \hline
\end{tabular}
\end{table}

Для настоящего исследования эти рамки рассматриваются как~\textbf{долгосрочный
пользовательский слой оценки}, а~не~как замена уже выполненных технических измерений.
Иными словами, текущий MVP по-прежнему опирается на~покрытие уровней
таксономии Блума, валидность структуры, поисково-фактологические
метрики и~операционную себестоимость. После коммерческого запуска
продуктовый контур оценки может быть усилен показателями
самоэффективности, ощущаемой компетентности, переноса обучения
в~практику и~достижения цели
~\cite{bandura_self_efficacy_2006, pcs_sdt_2026,
kirkpatrick_model_2026, shilts_retrospective_2008,
shogren_gas_2021}.

% ─────────────────────────────────────────────────────────────────
\section{Выбор и обоснование технологий}
% ─────────────────────────────────────────────────────────────────

\subsection{Фреймворк мультиагентной оркестрации}

Задача ИИ-методиста относится к классу процессов с~сохраняемым состоянием:
документ проходит этапы обработки, анализа источника, подготовки
структуры, наполнения черновыми материалами, проверки фактов, педагогической проверки
и~публикации.
Часть шагов выполняется последовательно, часть --- с условными ветками,
остановками и повторным запуском после решения с~участием человека.

В ходе отбора рассматривались LangGraph~1.x, AutoGen~0.4,
CrewAI~0.9, Agno~1.x и OpenAI Agents SDK
~\cite{langgraph_docs_2025, autogen_04_2024, crewai_docs_2025,
agno_docs_2025, openai_agents_sdk_2025}.

\begin{table}[h!]
  \caption{Сравнение фреймворков мультиагентной оркестрации}
  \label{tab:frameworks}
  \centering
  \begingroup
  \footnotesize
  \setlength{\tabcolsep}{4pt}
  \renewcommand{\arraystretch}{1.25}
  \begin{tabular}{|p{2.8cm}|c|c|c|c|p{2.2cm}|}
  \hline
  \textbf{Критерий}
    & \textbf{LangGraph}
    & \textbf{AutoGen}
    & \textbf{CrewAI}
    & \textbf{Agno}
    & \textbf{OpenAI Agents~SDK} \\
  \hline
  Модель исполнения
    & граф состояний
    & диалоговые агенты
    & ролевая команда
    & асинхронные агенты
    & агенты с~вызовом инструментов \\
  \hline
  Явный граф переходов
    & \textbf{+++}
    & ++
    & +
    & ++
    & + \\
  \hline
  Поддержка локальных LLM
    & \textbf{+}
    & \textbf{+}
    & \textbf{+}
    & \textbf{+}
    & --- \\
  \hline
  Контур участия человека
    & \textbf{+++}
    & ++
    & +
    & ++
    & + \\
  \hline
  Наблюдаемость и~трассировка
    & \textbf{+++}
    & ++
    & +
    & ++
    & ++ \\
  \hline
  Зрелость экосистемы
    & \textbf{+++}
    & ++
    & ++
    & +
    & ++ \\
  \hline
  \end{tabular}
  \endgroup
\end{table}

\textbf{Обоснование выбора LangGraph.}
LangGraph соответствует природе задачи: система должна хранить
состояние процесса, поддерживать возвраты к предыдущим шагам,
ручную проверку и трассировку каждого узла. Для этого \texttt{StateGraph} подходит лучше, чем
диалоговые модели. Дополнительным фактором стала интеграция
с LangSmith, которая даёт наблюдаемость на уровне отдельных вызовов
агентов и позволяет дебажить качество генерации.

\textbf{Причины отказа от альтернатив.}
AutoGen силён как среда многоагентного взаимодействия, но для
детерминированного конвейера требует дополнительных абстракций над
диалоговой моделью. CrewAI удобно описывает роли и задачи, однако
хуже соответствует продакшн-ориентированному асинхронному процессу.
Agno привлекателен своей асинхронностью и RAG-встроенностью, но
уступает в зрелости экосистемы. OpenAI Agents SDK не рассматривается
как базовый вариант из-за облачной привязки и отсутствия
траектории локального развёртывания.

\subsection{Выбранная архитектура}

После выбора оркестратора возникает второй вопрос: чем строить
контур обработки документов и поиска релевантных фрагментов.
В фактической MVP-реализации выбран интеграционный путь:
RAG-контур отвечает за индексирование и поиск по источнику,
граф знаний добавляет предметные связи и компетенции, а LangGraph
исполняет состояние агентного процесса поверх уже подготовленного
доказательного контекста.

Разделение соответствует сильным сторонам компонентов. RAG-контур
закрывает практическую задачу загрузки, индексации и поиска фрагментов
в контуре, приближённом к промышленному; граф знаний уменьшает зависимость от
плоского поиска по сходству; LangGraph хранит состояние курса, результаты
промежуточных нод и точки участия человека.

Итоговая техническая связка текущего MVP выглядит так:
Источник~$\rightarrow$ RAG-контур~$\rightarrow$ таблица свидетельств + граф знаний
$\rightarrow$ LangGraph~$\rightarrow$ артефакты курса. Это снижает связанность системы
и позволяет заменить поисковую подсистему без переписывания агентного графа.

Общая схема технологического стека приведена на рисунке~\ref{fig:tech-stack-flow}.

\begin{figure}[htbp]
  \centering
  \begin{tikzpicture}[
    node distance=1.2cm,
    every node/.style={font=\small},
    block/.style={rectangle, rounded corners=4pt, draw, minimum height=1cm, minimum width=2.4cm, align=center, fill=white},
    bigblock/.style={rectangle, rounded corners=4pt, draw, minimum height=1.2cm, minimum width=2.8cm, align=center, fill=blue!8, line width=0.8pt},
    arr/.style={-stealth, thick},
  ]
    % Row 1: Document -> Haystack AsyncPipeline -> SemanticSplitter -> OllamaEmbedder
    \node[block] (minio) {Document\\(MinIO)};
    \node[block, right=of minio] (pipeline) {Haystack\\AsyncPipeline};
    \node[block, right=of pipeline] (splitter) {Semantic\\Splitter};
    \node[block, right=of splitter] (embedder) {Ollama\\Embedder};

    % Row 2: Qdrant <- ... LangGraph StateGraph -> Course Artifacts
    \node[block, below=1.8cm of minio] (qdrant) {Qdrant};
    \node[bigblock, right=of qdrant] (langgraph) {LangGraph\\StateGraph};
    \node[block, right=of langgraph] (artifacts) {Course\\Artifacts};

    % Dashed box around Haystack components
    \begin{scope}[on background layer]
      \node[draw, dashed, rounded corners=6pt, inner sep=10pt,
            fit=(splitter)(embedder), label={[font=\footnotesize]above:Haystack 2.x}] {};
    \end{scope}

    % Row 1 arrows
    \draw[arr] (minio) -- (pipeline);
    \draw[arr] (pipeline) -- (splitter);
    \draw[arr] (splitter) -- (embedder);

    % Corner arrow from embedder down to qdrant
    \draw[arr] (embedder.south) -- ++(0,-0.6) -| (qdrant.north);

    % Row 2 arrows
    \draw[arr] (qdrant) -- (langgraph);
    \draw[arr] (langgraph) -- (artifacts);
  \end{tikzpicture}
  \caption{Технологический стек ИИ-методиста: контур обработки и генерации}
  \label{fig:tech-stack-flow}
\end{figure}


\subsection{Большая языковая модель}

На этапе экспериментального MVP сравнивались три направления:
облачная модель общего назначения, российский облачный провайдер
и локальная открытая модель.

\begin{table}[h!]
  \caption{Сравнение LLM-вариантов для текущей итерации}
  \label{tab:llms}
  \centering
  \begin{tabular}{|p{4cm}|p{3.2cm}|p{3.2cm}|p{3.2cm}|}
  \hline
  \textbf{Критерий}
    & \textbf{GPT-4o}
    & \textbf{YandexGPT 4}
    & \textbf{Qwen2.5 7B} \\
  \hline
  Способ доступа
    & облачный API
    & облачный API
    & \textbf{локально через Ollama} \\
  \hline
  Суверенитет данных
    & ---
    & $\pm$
    & \textbf{+} \\
  \hline
  Стоимость итераций и отладки
    & высокая
    & средняя
    & \textbf{низкая} \\
  \hline
  Русскоязычный контент
    & высокая
    & высокая
    & достаточная для базовой конфигурации \\
  \hline
  Инфраструктурная простота
    & высокая
    & высокая
    & \textbf{высокая для локального стенда} \\
  \hline
  \end{tabular}
\end{table}

Для текущих экспериментов базовой конфигурации выбрана модель
\texttt{qwen2.5:7b} через Ollama. Причина выбора прагматична:
эта модель доступна локально, даёт воспроизводимые результаты,
не требует внешнего API и позволяет многократно прогонять
эксперименты на ядровом корпусе без роста переменных затрат.

Это \textbf{не окончательное ограничение} архитектуры. ИИ-методист проектируется как провайдер-агностичная
система: при необходимости локальная базовая конфигурация может быть заменена
на более крупный вариант с открытыми весами семейства Qwen~2.5 или
на облачный провайдер, если сценарий внедрения это допускает.

\subsection{Поиск, векторные представления и граф знаний}

Для поискового слоя в текущей итерации зафиксирована не конкретная
векторная подсистема, а контракт поискового слоя: на входе источник,
на выходе нумерованные фрагменты таблицы свидетельств, пригодные для
цитирования и проверки фактов. В реализации этот контракт закрывается
RAG-контуром и графом знаний. В production-коде RAG подключён через
порт \texttt{DifyRAGApi} и retrieval-адаптер к Dify API: ключ и датасет
выбираются по tenant-контексту, базовый метод поиска ---
\texttt{keyword\_search}, \texttt{top\_k=5}, порог релевантности
\texttt{score\_threshold=0.5}, reranking отключён. На выходе формируются
ранжированные фрагменты \texttt{RAGChunk}: текст, оценка релевантности,
имя документа, позиция в источнике и ключевые слова.
Отдельный платформенный индекс методологических знаний в контуре
\texttt{ai-orchestration} использует Qdrant-коллекцию
\texttt{methodology\_knowledge\_v2}, hash-эмбеддинги, размерность
вектора 64 и тайм-аут запроса 10~с; он рассматривается как
совместимый инфраструктурный слой, а не как единственное доказательство
работоспособности MVP.

Выбор такого контракта объясняется четырьмя причинами:

\begin{itemize}
  \item возможность сохранять привязку к источнику независимо от конкретного
    векторного хранилища;
  \item поддержка локального развёртывания и запрета внешней передачи
    документов;
  \item возможность обогащать найденные фрагменты предметными связями
    графа знаний;
  \item совместимость с заменой retrieval-адаптера или подключением
    Qdrant/локальных векторных представлений без изменения узла LangGraph.
\end{itemize}

Для поискового слоя выполнена детерминированная проверка по схеме RAG
на 30 вопросно-ответных парах: достоверность ответа составила 0.335,
релевантность ответа --- 0.239, точность контекста --- 0.567,
полнота контекста --- 0.335. Эти значения
используются как диагностический слой: текущий контур оценки RAG работает
как консервативный контроль перед проверкой человеком, то есть помогает раньше
отправить рискованные фрагменты на проверку человеком. Они не доказывают
качество поиска; качество поиска требует расширения корпуса и
экспертной разметки релевантных фрагментов. Основной продуктовый
показатель H2 при этом считается производной метрикой привязки к источнику:
долей фактических утверждений, подтверждённых таблицей свидетельств в финальном
\texttt{verification\_report}.

\subsection{Стратегии разбиения документов на фрагменты}

Для обработки документов были рассмотрены две стратегии:
рекурсивное фиксированное разбиение и семантическое разбиение.
Детерминированная проверка по схеме RAG дала точность контекста 0.567 и
полноту контекста 0.335; поэтому текущий RAG-контур используется как
рабочий контрольный слой для MVP, но не считается доказанным по качеству
поиска. Архитектурно более перспективным остаётся семантическое
разбиение внутри поисковой подсистемы и последующая экспертная оценка
релевантности найденных фрагментов.
Конкретная реализация может быть заменена: в текущем MVP это
RAG-контур с Dify retrieval-адаптером, в долгосрочной инфраструктурной
версии допустим переход на Haystack/Qdrant
при сохранении контракта таблицы свидетельств.

Учебные и нормативные документы содержат смысловые блоки разной длины: определения, процедуры,
ограничения, списки требований. Жёсткое разбиение по фиксированному размеру разрушает
эти границы и создаёт более шумный поисковый контекст. Семантическое
разбиение лучше сохраняет логические единицы, что особенно важно
для нормативных документов и продуктово-редакционных руководств.

\subsection{Перспективное развитие генерации изображений}

Для мультимедийного слоя сравнивались DALL-E~3 и FLUX.1-dev
~\cite{dalle3_2023, flux_2024}. DALL-E~3 удобен как внешний сервис
с высоким качеством и быстрым стартом, но вносит в систему
API-зависимость и переменные затраты. FLUX.1-dev допускает локальное
развёртывание и потому лучше соответствует логике продукта,
ориентированного на локальный контур.

В результате для перспективного слоя развития выбрана следующая стратегия:

\begin{itemize}
  \item \textbf{по умолчанию} --- локальный FLUX.1-dev для черновых
    и полуфинальных образовательных иллюстраций;
  \item \textbf{опционально} --- внешний провайдер класса DALL-E~3
    для сценариев повышенного качества, где требуется более «маркетинговое»
    качество изображения.
\end{itemize}

Решение не закрывает дорогу к внешнему качественному генератору,
но сохраняет базовый маршрут локального развёртывания для большинства
B2B-сценариев.

\subsection[Перспективное развитие генерации подкастов]{Перспективное развитие генерации подкастов на основе синтеза речи}

Для генерации русскоязычных подкастов сравнивались Silero~TTS~v4
и ElevenLabs Multilingual~v2~\cite{silero_tts_2024, elevenlabs_2024}.
На текущем этапе различие между ними определяется прежде всего
не субъективным MOS, а эксплуатационной логикой.

\begin{itemize}
  \item \textbf{Silero~TTS} подходит для локального развёртывания,
    не требует внешнего API и по заложенной в оценочную модель
    оценке даёт стоимость порядка 4~руб.\ на курс за счёт только
    GPU-времени.
  \item \textbf{ElevenLabs} обеспечивает более «человеческую»
    интонационную подачу, но при оценке 40{,}000 символов на курс
    стоимость составляет порядка 1{,}080~руб.\ на курс, что делает
    его скорее опцией повышенного качества, чем базовым слоем синтеза речи для MVP.
\end{itemize}

Панельная MOS-оценка пока не завершена. Для этого уже подготовлен
отдельный оценочный скрипт, однако в текущей итерации для выбора
технологии решающими оказались суверенитет данных и предсказуемость
себестоимости. Поэтому для будущего мультимедийного слоя базовым
движком синтеза речи выбран Silero, но этот слой не используется как
доказательство текущей готовности MVP.

% ─────────────────────────────────────────────────────────────────
\section{Проверенные результаты экспериментов}
\label{sec:exp-results}
% ─────────────────────────────────────────────────────────────────

\subsection[Результаты одношаговой LLM-конфигурации (B2)]{Измеренные результаты одношаговой LLM-конфигурации (B2)}

Первым полностью воспроизводимым экспериментом стал лабораторный прогон
одношаговой LLM-конфигурации (B2) на~измеренном 19-документном корпусе:
продуктово-редакционные документы, корпоративные регламенты и учебные
программы MIT~OCW. Эти числа используются как технический базовый ориентир
текущего MVP и синхронизированы с артефактом базового сравнения.
В этой конфигурации использовалась локальная модель \texttt{qwen2.5:7b},
которая по одному запросу генерировала структуру курса в JSON-формате
без поиска по корпусу, без проверки фактов и без отдельных специализированных
агентов.

\begin{table}[h!]
  \caption{Измеренные результаты одношаговой LLM-конфигурации (B2)}
  \label{tab:b2prelim}
  \centering
  {\small
  \setlength{\tabcolsep}{3pt}
  \renewcommand{\arraystretch}{1.15}
  \begin{tabular}{|p{3.8cm}|p{0.8cm}|p{2.0cm}|p{2.1cm}|p{2.4cm}|p{3.0cm}|}
  \hline
  \textbf{Срез корпуса}
    & \textbf{$n$}
    & \textbf{Среднее время, с}
    & \textbf{Среднее покрытие, \%}
    & \textbf{Доля курсов $\geq 70\%$}
    & \textbf{Структурная валидность} \\
  \hline
  Все документы
    & 19
    & 4.37
    & 61.4
    & 21.1\%
    & 100\% \\
  \hline
  Продуктово-редакционные документы
    & 7
    & 5.65
    & 73.8
    & 28.6\%
    & 100\% \\
  \hline
  Корпоративные регламенты
    & 6
    & 3.65
    & 38.9
    & 0\%
    & 100\% \\
  \hline
  Учебные программы MIT~OCW
    & 6
    & 3.61
    & 69.5
    & 33.3\%
    & 100\% \\
  \hline
  \end{tabular}
  }
\end{table}

Распределение покрытия по категориям визуализировано на рисунке~\ref{fig:bloom-by-category}.

\begin{figure}[htbp]
\centering
\begin{tikzpicture}
\begin{axis}[
    ybar,
    width=0.9\textwidth,
    height=7cm,
    ylabel={Покрытие уровней Блума, \%},
    ymin=0, ymax=100,
    xtick=data,
    symbolic x coords={Нормативные,{Все документы},{Учебные программы},Продуктовые},
    x tick label style={font=\small, rotate=25, anchor=east},
    bar width=22pt,
    enlarge x limits=0.2,
    nodes near coords,
    nodes near coords style={font=\small, anchor=south},
]
\addplot[fill=red!50, draw=red!70!black] coordinates
    {(Нормативные,40) ({Все документы},61.5) ({Учебные программы},66.7) (Продуктовые,75)};

% Целевой порог
\draw[dashed, thick, black!70] (axis cs:Нормативные,70) -- (axis cs:Продуктовые,70)
    node[pos=0.75, above, font=\scriptsize] {Целевой порог (70\,\%)};

% Перекрасить Продуктовые (>=70) — наложим поверх
\addplot[fill=teal!70!black, draw=teal!50!black, forget plot] coordinates
    {(Нормативные,0) ({Все документы},0) ({Учебные программы},0) (Продуктовые,75)};
\end{axis}
\end{tikzpicture}
\caption{Покрытие уровней таксономии Блума по категориям документов (конфигурация B2)}
\label{fig:bloom-by-category}
\end{figure}


Частота появления отдельных уровней таксономии Блума показана на рисунке~\ref{fig:bloom-levels}.

\begin{figure}[htbp]
\centering
\begin{tikzpicture}
\begin{axis}[
    compat=1.18,
    xbar,
    width=0.85\textwidth,
    height=7cm,
    xlabel={Доля документов, \%},
    xmin=0, xmax=110,
    ytick=data,
    yticklabels={
        {Создание (Create)},
        {Оценка (Evaluate)},
        {Анализ (Analyze)},
        {Применение (Apply)},
        {Понимание (Understand)},
        {Запоминание (Remember)}
    },
    y tick label style={font=\small, anchor=east},
    nodes near coords,
    nodes near coords align={horizontal},
    nodes near coords style={font=\small, /pgf/number format/.cd, fixed, precision=1},
    bar width=12pt,
    enlarge y limits=0.15,
    colormap={bloom}{color(0)=(teal!30!white) color(1)=(teal!90!black)},
    colormap access=map,
    point meta=explicit,
]
\addplot coordinates {
    (12.5,0)  [0.0]
    (18.75,1) [0.2]
    (37.5,2)  [0.4]
    (62.5,3)  [0.6]
    (87.5,4)  [0.8]
    (100,5)   [1.0]
};
\end{axis}
\end{tikzpicture}
\caption{Частота покрытия уровней таксономии Блума в корпусе (n=16, конфигурация B2)}
\label{fig:bloom-levels}
\end{figure}


Полученные данные показывают, что одношаговая конфигурация обладает двумя
сильными сторонами: она очень быстро генерирует структурно валидный
ответ и не требует сложной инфраструктуры. Однако эти преимущества
не решают ключевую задачу исследования.

Во-первых, среднее покрытие уровней таксономии Блума составляет лишь 61.4\,\%,
а порог 70\,\% достигается только в 21.1\,\% случаев. Во-вторых,
базовая конфигурация значительно лучше справляется с продуктово-редакционными
документами, чем с нормативными: для официальных нормативных документов покрытие падает
до 38.9\,\%, что означает смещение к нижним когнитивным уровням
запоминания и понимания. В-третьих, плотность ссылок на источник равна нулю:
модель не даёт никакой явной привязки результата к источнику,
поэтому одношаговая LLM-конфигурация (B2) не может считаться достаточной для
нормативных и знаниево-критичных сценариев.

Время генерации по категориям документов представлено на рисунке~\ref{fig:generation-time}.

\begin{figure}[htbp]
  \centering
  \begin{tikzpicture}
    \begin{axis}[
      ybar,
      width=0.85\textwidth,
      height=7cm,
      ylabel={Время, с},
      symbolic x coords={Продуктовые, Регламенты, {MIT OCW}, {Все документы (среднее)}},
      xtick=data,
      x tick label style={font=\small, rotate=30, anchor=east},
      ymin=0, ymax=7,
      bar width=1.2cm,
      nodes near coords,
      nodes near coords align={vertical},
      every node near coord/.append style={font=\small},
      enlarge x limits=0.2,
      grid=major,
      ymajorgrids=true,
      xmajorgrids=false,
    ]
      \addplot coordinates {
        (Продуктовые, 5.65)
        (Регламенты, 3.65)
        ({MIT OCW}, 3.61)
        ({Все документы (среднее)}, 4.37)
      };
    \end{axis}
  \end{tikzpicture}
  \caption{Среднее время генерации структуры курса по категориям для одношаговой LLM-конфигурации (B2)}
  \label{fig:generation-time}
\end{figure}


Именно этот набор ограничений и мотивирует переход от одношаговой LLM (B2)
к ИИ-методисту (B4):
необходима архитектура, которая сохраняет скорость, но усиливает
поиск по корпусу, проверка фактов и глубину методической структуры.

\subsection[Сравнение LLM-конфигурации (B2) и ИИ-методиста (B4)]{Сравнение одношаговой LLM (B2) и ИИ-методиста (B4) на одном корпусе}

После стабилизации REST-контура ИИ-методиста выполнен парный прогон
одношаговой LLM (B2) и ИИ-методиста (B4) на том же 19-документном корпусе.
Для каждого документа
создавался отдельный \texttt{kb\_slug} B4 формата
\texttt{diploma-b4-001}, в~запуск передавался ровно один Markdown-файл,
а результат сопоставлялся с ранее измеренной строкой B2 для того же
\texttt{doc\_id}. Это устраняет риск некорректного сравнения средних
из разных прогонов: итоговая метрика считается как
$\Delta_i = B4_i - B2_i$ для каждого документа.

\begin{table}[h!]
  \caption{Сравнение одношаговой LLM (B2) и ИИ-методиста (B4) на 19 документах}
  \label{tab:paired-b2-b4}
  \centering
  \begin{tabular}{|p{5.8cm}|p{2.5cm}|p{4.8cm}|}
  \hline
  \textbf{Метрика} & \textbf{Значение} & \textbf{Вывод} \\
  \hline
  Среднее покрытие Блума B2 & 61.4\,\% & базовая конфигурация ниже целевого порога \\
  \hline
  Среднее покрытие Блума B4 & 83.3\,\% & B4 проходит порог $\geq 70\,\%$ \\
  \hline
  Средняя парная дельта Блума & +21.9 п.п. & прирост считается по тем же документам \\
  \hline
  Документы с положительной дельтой & 15 из 19 & B4 лучше B2 на большинстве корпуса \\
  \hline
  Привязка к источнику B4 & 100\,\% & B4 даёт проверяемую опору на источник \\
  \hline
  Знак качества B4 & 100\,\% «пройдено» & все B4-запуски готовы к проверке человеком \\
  \hline
  \end{tabular}
\end{table}

Наиболее сильный эффект наблюдается на корпоративных регламентах:
среднее покрытие Блума увеличилось с 38.9\,\% до 83.3\,\%,
то есть на +44.4 п.п. Для продуктово-редакционных документов
дельта составила +9.5 п.п., для программ MIT~OCW --- +13.8 п.п.
Таким образом, B4 закрывает именно тот класс документов, где B2 был
наиболее слаб: нормативные и знаниево-критичные материалы, требующие
явной привязки к источнику и проверки качества.

\subsection{Оценочные наборы для RAG и проверки фактов}

Парный прогон одношаговой LLM (B2) и ИИ-методиста (B4) закрывает
основной технический разрыв между базовой конфигурацией и ИИ-методистом, однако для более тонкой оценки RAG-компонента
и проверки фактов использованы дополнительные наборы.

\begin{table}[h!]
  \caption{Оценочные наборы текущего MVP}
  \label{tab:evalassets}
  \centering
  \begin{tabular}{|p{5.5cm}|p{2.5cm}|p{6.0cm}|}
  \hline
  \textbf{Артефакт}
    & \textbf{Объём}
    & \textbf{Назначение} \\
  \hline
  Измеренный корпус одношаговой LLM (B2)
    & 19
    & базовая и конвейерная подготовка курсов \\
  \hline
  Вопросно-ответные пары (\texttt{qa\_pairs.json})
    & 30
    & оценка по схеме RAG: достоверность ответа, релевантность ответа,
      точность контекста и полнота контекста \\
  \hline
  Фактуальные утверждения (\texttt{statements.json})
    & 100
    & оценка точности, полноты и F1-меры для~проверки фактов \\
  \hline
  \end{tabular}
\end{table}

Артефакты превращают работу из описательной архитектурной записки
в экспериментально воспроизводимую постановку. На этих наборах выполнен
детерминированный прогон лабораторной оценки: для оценки по схеме RAG получены
достоверность ответа 0.335, релевантность ответа 0.239, точность контекста 0.567
и полнота контекста 0.335; для проверки фактов получены точность 0.600,
полнота 0.900 и F1-мера 0.720. Высокая полнота выбрана сознательно: для
методического ассистента опаснее пропустить неподтверждённое утверждение,
чем отправить лишний фрагмент на ручную проверку.
Следовательно, текущая лабораторная оценка подтверждает пригодность
связки проверки фактов и RAG как консервативного контрольного слоя перед
проверкой человеком, но не является доказательством полного качества поиска.
Для такого вывода требуется расширенный корпус, экспертная разметка
релевантных контекстов и повторение замеров на открытых курсах разных
жанров.

\subsection[Сценарная экономика единицы продукта MVP]{Сценарная экономика единицы продукта MVP на основе проверенного прогона}

Помимо технических метрик, собрана сценарная оценка
экономики единицы продукта. Речь идёт не о подтверждённой рыночной выручке,
а о расчётной модели, основанной на текущей конфигурации конвейера
и внутренних допущениях о цене привлечения и контрактной монетизации.

\begin{table}[h!]
  \caption{Сценарная экономика единицы продукта текущего MVP на основе проверенного прогона}
  \label{tab:unitecon}
  \centering
  \begin{tabular}{|p{6.0cm}|p{3.0cm}|p{4.5cm}|}
  \hline
  \textbf{Показатель}
    & \textbf{Значение}
    & \textbf{Интерпретация} \\
  \hline
  Прямая себестоимость генерации курса
    & 75 руб.
    & инфраструктурная себестоимость курса в текущем конвейере \\
  \hline
  Отношение к стоимости подрядной разработки
    & 0.05\%
    & текущий MVP в модели существенно дешевле внешнего производства \\
  \hline
  LTV / CAC
    & 10.0
    & при выбранных допущениях модель не выглядит заведомо убыточной \\
  \hline
  Срок окупаемости
    & 2 мес.
    & возврат затрат возможен в коротком цикле продаж \\
  \hline
  Порог безубыточности
    & 14 клиентов
    & ориентир для коммерческой стадии после MVP \\
  \hline
  \end{tabular}
\end{table}

Цифры служат проверкой того, что выбранная архитектура
не разрушает экономику продукта: после завершения парного прогона B4
система остаётся технически дешёвой в эксплуатации.

Одновременно эти значения не должны интерпретироваться как окончательное
доказательство соответствия продукта рынку. В рамках ВКР коммерческий блок является
сценарной моделью устойчивости, а не утверждением о спросе, выручке или
проверке на внешних закрытых данных.

\subsection[Ограничения этапа и план исследования]{Ограничения текущего этапа и план абляционного исследования}

Наиболее существенное ограничение текущего этапа состоит уже не в
отсутствии B4-прогона: парное B2/B4-сравнение на 19 документах выполнено
и сохранено в \texttt{paired\_b2\_b4\_experiment.json}. Оставшееся
ограничение связано с тем, что автоматические метрики ещё не заменяют
слепую экспертную оценку и абляционное исследование вклада отдельных
узлов конвейера.

Ограничение накладывает границу на интерпретацию результатов:

\begin{itemize}
  \item текущие цифры главы~2 следует трактовать как
    \textbf{автоматическую парную проверку одношаговой LLM (B2)
    и ИИ-методиста (B4)}, а не как замену
    экспертной педагогической экспертизы;
  \item гипотезы H1--H3 в данной работе проверяются в пределах
    внутреннего / открытого корпуса, парной проверки B2/B4 и лабораторной оценки;
  \item $n = 19$ является лабораторной проверкой MVP, но недостаточен для
    сильного вывода о превосходстве над аналогами; такой вывод
    должен проверяться отдельной сравнительной проверкой на открытых курсах $n \geq 30$;
  \item сравнение с ручной разработкой (B1), AI-конструктором (B3)
    и слепая экспертная педагогическая экспертиза
    остаются полезным расширением методики, но не входят в заявленный
    объём доказательств текущей ВКР;
  \item абляционное исследование является направлением развития
    доказательной базы, так как оно позволит отдельно измерить вклад
    поискового контура, проверки фактов и агентов проверки в итоговое качество.
\end{itemize}

Глава~2 фиксирует и достигнутые результаты, и текущую границу
зрелости системы, определяя точку отсчёта для дальнейших экспериментов.

\subsection{Финальная MVP-конфигурация}

На основании проведённого отбора и измеренных результатов
итоговая MVP-конфигурация ИИ-методиста строится на следующей
конфигурации:

\begin{itemize}
  \item \textbf{Локальное LLM-ядро}: Qwen через OpenAI-совместимую конечную точку Ollama;
  \item \textbf{Оркестратор}: LangGraph~1.x (\texttt{StateGraph});
  \item \textbf{Контур поиска}: RAG-инструмент с параметрами
    \texttt{keyword\_search}, \texttt{top\_k=5},
    \texttt{score\_threshold=0.5} и таблица свидетельств;
  \item \textbf{Семантическое усиление}: граф знаний;
  \item \textbf{Состояние и результаты}: PostgreSQL-хранилище контрольных точек и результатов;
  \item \textbf{Контроль качества}: проверка фактов + верификация + знак качества;
  \item \textbf{План развития}: Qdrant/локальные векторные представления, FLUX, Silero,
    SAM и Backward~Design после парной проверки основного B4-контура.
\end{itemize}

Конфигурация отвечает целям работы: локальное развёртывание,
воспроизводимые эксперименты, приемлемая себестоимость, усиление
поискового слоя и педагогической структуры относительно одношаговой
базовой конфигурации.
